\documentclass[10pt]{article}

\usepackage[utf8]{inputenc}

\usepackage[T1]{fontenc}

\usepackage{amsmath}

\usepackage{amsfonts}

\usepackage{amssymb}

\usepackage[version=4]{mhchem}

\usepackage{stmaryrd}

\usepackage{graphicx}

\usepackage[export]{adjustbox}

\graphicspath{ {./images/} }

\usepackage{mathrsfs}



\begin{document}

Непосредственным обобщением коники в $P_{3}$ является квадратичныї әлечент - $(n-2)$-мерная невырожденная квадрика в $p_{n}(n>3)$. Многообразием $(h, m, n)^{2}$ в пространстве $P_{n}$ наз. $m$-параметрич. семейство квадратичных әлементов, гиперплоскости к-рых образуют $h$-параметрич. семейство. Многообразия $(h, h, 3)^{2}, h=1$, 2,3 , являются наиболее общими однопараметрич. семействами, конгруэнциями и комплексами коник в $P_{3}$. С каждым локальным квадратичным элементом многообразия $(h, h, n)^{2}$ шри $h<n$ ассоциируется $(n-h-1)-$ мерное характеристич. подпространство и $(h-1)$-мерное полярное нодпространство. Р а н г о м м н о гоо б р а 3 и я $(h, h, n)^{2}$ наз. число $R$, равное $n-h-1-v$, где $v$ - размерность подпространства, по к-рому характеристич. подпространство пересекается с его полярным пространством.



Дифференциальную геометриго многообразия $(n, n$, $n)^{2}$ можно рассматривать как геометрию нек-рой регулярной гиперповерхности $(n+1)$-мерного тангенцального центропроективного пространства $P_{0}^{n+1}$, в к-ром исходное $n$-мерное пространство играет роль неиодвижной точки.



Квадрика размерности $p \leqslant n-2$ пространства $P_{n}$ всегда лежит в $(p+1)$-мерном подиространстве. Алгебраич. поверхности порядка $k>2$ әтим свойством в общем случае не обладают. При исследовании многообразий алгебраит. поверхностей оказалось целесообразным выделить семейства говерхностей таких, к-рые лежат в плоскостя х на единицу большей размерности. Невырожденная $(n-2)$-мервая алгебраич. поверхность порядка $k$, принадлежащая гиперплоскости $P_{n-1}$ простравства $P_{n}, \quad$ наз. пл о ским а лге б ра и ч. элем ентом порядка $k$. Многообразием $(h, m, n)^{k}$ наз. $m$-мерное многообразие плоских алгебраич. элементов порядка $k$, гиперплоскости к-рых образуют $h$-параметрич. семейство. Фундаментальный объект 1-го порядка является основным объектом мвогообразия $(h, m, n)^{k}$.



В применениях геометрии к гидромеханике, теории поля и дифференциальным уравнениям используются многообразия линий и поверхностей, не являгощихся в общем случае алгебраическими. Изучение таких многообразий представляет и самостоятельный интерес для геометрии. $К \mathbf{p}$ и в ли нейн о й кон г ру нЦ и е й в трехмерном пространстве наз. двупараметрич. семейство $\Gamma_{t}$ кривых $x=x(t, u, v)$ такое, что через каждую тотку пространства проходит в общем случае единственная кривая семейства. При фиксированных $u$ и $v$ выделяется кривая $C_{t}$ семейства $\Gamma_{t} ;$ пи фиксированном $t$ и переменных $u, v$ выделяется поверхность $s_{u v}$, называемая т р а с в е р с а льн о ӥ П о в рх нос тью конгруэнции $\Gamma_{t}$. Точки кривой $C_{t}$, в к-рых $\left(x x_{t} x_{u} x_{v}\right)=0$, наз. фо к а льны м и то ч к а ми. Совокупность фокальных тотек кривых $C_{t}$ конгруэнции $\Gamma_{t}$ образует ф о к а л ь у у ю П о в е р х н о т ь конгруәнции. Поверхность $v=v(u)$ конгруэнтии $\Gamma_{t}$, на к-рой кривые $C_{t}$ имеют огибающую, наз. І ІІ а в но й п о в е р х н о тью конгруәнции $\Gamma_{t}$.



Лum.: [1] М а л а х о в с к и й В. С., в кн.: Итогъ науки и техники. Проблемы геометрии, т. 12 , м., 1981 , с. $31-60$.



ФКГУРА - подмножество $F$ однородного пространства $E_{n}$ с фундаментальной группой $G$, к-рое можно включить в систему $R(F)$ шодмножеств әтого пространства, изоморфнуіо нек-рому пространству геометрич. объекта Ф (см. Іеометрических объектов теория). Множество $R(F)$ наз. І р о с т р а н с т в о м ф и г ур ы $F$. Компоненты геометрич. объекта Ф наз. к о о рд и н а т а м и соответствугщей $\Phi$. $F$. Каждой $F$ пространства $E_{n}$ соответствует класс $\{\Phi\}$ подобных геометрич. объектов. Ранг, жанр, характеристика и тип геометрич. объекта Ф класса $\{\Phi\}$ наз. р а н го м, ж а н р о, х а р а т теристи ко й и типом Ф) и г у р ы $F$ (т.н. а риф ме тиче с ки е ин в ар и а н т ы ф и г у р ы). Напр., окружность в трехмерном евклидовом пространстве является фигурой ранга 6, жанра 1, характеристики 1, тиша 1; точка в трехмерном проективном пространстве - фигурой ранга 3, жанра 0 , характеристики 2, типа 1. Вполне интегрируемая система уравнениӥ Пфофффа, определяющая геометрический объект Ф, наз.. с и с т е-



\includegraphics[max width=\textwidth, center]{2023_07_09_c7449e77e62a0f371846g-1}

г у р ы $F$.



Пусть $F$ и $\bar{F}$ - две $\Phi$. Пространства $E_{n}$. Если существует отобраэкение п ространства $R(F)$ на пространство $R(\vec{F})$, при к-ром любой геометрич. объект, соответствующий Ф. $\bar{F}$, охватывается любым геометрич. объектом, соответствующим Ф. $F$, то говорят, что Ф. $F$ о х в аты в а т или инд у и и у е т Ф. $\bar{F}$ (Ф. $\bar{F}$ охватывается или индуцируется Ф. $F$ ). Ф. $F$ ранга $N$ наз. п р о с т о й. если она не охватывает никакой другой Ф. меньшего ранга. Ф. $F$ наз. и н д у ц и р у ю и е й ф иг у р ой и нд е к с а $\bar{N}<N$, если существует охватываемая ею Ф. ранга $\bar{N}$, причем ранг $N^{\prime}$ любой другой Ф. $F^{\prime}$, охватываемой Ф. $F$, не превосходит $\bar{N}$. Напр., точка, p-мерная плоскость, гиперквадрика в $n$-мерном проективном пространстве являются простыми Ф., a гиперквадрика в $n$-мерном аффинном пространстве и $d$-мерная $(d \leqslant n-2)$ квадрика в $n$-мерном проективном пространстве являются индуцирующими Ф. соответственно индексов $n$ и $(d+2)(n-d-1)$.



П а р ой фи г у р $F=\left(F_{1}, F_{2}\right)$ наз. упорядоченная совокупность двух Ф. Ко бфф иц иен том инцид е н т о с т и пары Ф. наз. число $k=N_{1}+N_{2}-N$, где $N_{i}(i=1,2)$ ранг фигуры $F_{i}$, а $N$-ранг системы форм $\Omega_{1}, \Omega_{2}, J_{i}=1,2, \ldots, N_{i}$, являючцхся левыми частями уравнений стационарности Ф. $F_{1}$ и $F_{2}$. Если $k=0$, то пара $F=\left(F_{1}, F_{2}\right)$ наз. н е и н ци п $е$ н т но й.



Лum.: [1] JI а п т е в Г. Ф., “Тр. Моск. матем. об-ва", 1953, т. 2, с. $275-383$; [2] М а л а х о в с к и й В. С., «Тр. Геометр. семинара. Ин-т научн. информ. АН СССР», 1969, т. 2, с. $179-$ 206 : [3] е г о ж е, в сб.: Итоги науки и техники. Алгебра. То-

пология. Геометрия, т. 10, М., 1972, с. $113-58$.



ФИГУРНОЕ ЧИСЈО - см. А рифметический ряд. ФИДУ ЦАЛЬНОЕ РАСПРЕДЕЛЕНИЕ-распределение $P_{x}^{*}$ параметра $\theta$ семейства распределений $\mathscr{P}=$ $=\left\{P_{\theta}, \quad \theta \in \Theta\right\}$ наблюдения $x$. Введено Р. Фишером [1] для числовых $\theta$ и $x$ в случае, когда функция распределения $F(x ; \theta)$ наблюдения $x$ убывает с ростом $\theta$ так, что $F^{*}(\theta \mid x)=1-F(x \mid \theta)$, рассматриваемая как функция от $\theta$ при фиксированном $x$, обладает свойствами функции распределения (к такой ситуации часто приводит использование в качестве $x$ достаточной статистикИ).



Ф. р. определено для инвариантных семейств распределений (см. [2]-(4]). Именно, пусть группа $G$ преобразований $g$ действует на множествах $X$ и $\Theta$. Семейство распределений наз. и н в а р и а н т н ы м, если $g x$ имеет распределение $P_{g \theta}$, когда $x$ имеет распределение $P_{\theta}$. В этом случае рассматривают эквивариантные регпающие правила $\delta: X \longrightarrow D$ (такие, тто $\delta(g x)=g \delta(x)$ для всех $x$ и $g)$ и инвариантные функции потерь $L_{\theta}(d)$ (такие, что $L_{g \theta}(g d)=L_{\theta}(d)$ для всех $\theta$, $d$ и $g)$. Если действие группы $G$ на множестве $\Theta$ транзитивно, то семейство $\mathscr{J}$ обладает определенным свойством однородности: для фикси рованного значения параметра $\theta_{0}$ и наблюдения $x$, имеющего распределение $P_{\theta_{0}}$, распределение $g x$ пробегает все семейство $\mathfrak{g}$, когда $g$ II робегает $G$. Пусть $D$-множество вероятностных мер на $\Theta$ (предполагаются заданными нек-рые б-алгебры $\mathscr{B}(\Theta)$ и $\mathscr{B}(X)$, так что преобразования группы $G$ измеримы). Пусть действие $G$ на $D$ задано соотнопе-





\end{document}